\emph{The seeding requirement.} By Proposition \claimref{dark-matter.seeding-gravity-floor-dark-towers-temperature} a tower heated by graviton-mediated production at rate $c\,T^{8}/M_{P}^{4}$ (Fact \claimref{dark-matter.imported-massless-graviton-production-rate}) carries comoving yield $Y = c\,T_{\mathrm{desc}}^{3}/(3\kappa_{s}M_{P}^{3})$, $\kappa_{s} = 145.3$ at $g_{*} = g_{s} = 100$, $M_{P} = 2.435 \times 10^{18}$ GeV. Summing over the produced towers, $c$ is the summed coefficient over all produced light species. The provenance of the rate absorbs the particle-pair and mean-energy factors into at most $2^{1/4}$ on the temperature ratio $\xi \propto Y^{1/4}$, so at most a factor $2$ on $Y$; we allow $Y \le 2c\,T_{\mathrm{desc}}^{3}/(3\kappa_{s}M_{P}^{3})$. If the relics' mass is $m \le T_{\mathrm{desc}}$ (the bath condition of T14's signpost), their mass per unit entropy is at most $2c\,T_{\mathrm{desc}}^{4}/(3\kappa_{s}M_{P}^{3})$. The remainder is $(\Omega_{c}/\Omega_{b} - 1)$ times the baryon mass per unit entropy, $m_{p}\,\eta_{s}$ with $\eta_{s} = 8.7 \times 10^{-11}$ the baryons per unit entropy of Proposition \claimref{dark-matter.coefficients-floors-they-give-what-gravity}(iii) and $\Omega_{c}/\Omega_{b} - 1 = 4.364 \pm 0.065$ from Proposition \claimref{dark-matter.population-count-one-share-half-percent}. Taking the lower ends, $4.299$ and $8.65 \times 10^{-11}$, and $m_{p} = 0.938272$ GeV, the requirement is $c\,T_{\mathrm{desc}}^{4} \ge A/2$ with
\[
A = 3\kappa_{s}M_{P}^{3}\,(4.299)(0.938272)(8.65 \times 10^{-11})\ \mathrm{GeV}, \qquad (1.217 \times 10^{12})^{4} < A < (1.218 \times 10^{12})^{4}\ \mathrm{GeV}^{4}.
\]
(With central inputs and no factor $2$ this is T14's signpost: $T_{\mathrm{desc}} \ge 1.22 \times 10^{12}\,c^{-1/4}$ GeV, above $2 \times 10^{12}$ GeV for $c \le 0.14$.)

\emph{Premise (R).} Premise \claimref{neutrino.r-descent-reheats-below-lower-heavy} states $T_{\mathrm{desc}} \le \alpha = \min(\alpha, M_{3})$, hence $T_{\mathrm{desc}} \le M_{3}$, the hypothesis of the statement. We use only this; the branch value $\alpha \approx M_{3}/66$ of T15 is a fitted instance and would strengthen (i) and (ii) by $66^{4}$ in $c^{*}$.

\emph{(i)} Proposition \claimref{standard-model.what-angle-asks-neutrino-sector} gives, on the chain with the $(15,2,2)$ at central values, $v_{R} = 4.2 \times 10^{10}$ GeV and $m_{D}^{2} = m_{3}f_{J}v_{R}$; the seesaw $m_{3} = m_{D}^{2}/M_{3}$ of Proposition \claimref{dark-matter.ceiling-valued-inherited-success} then reads $M_{3} = f_{J}v_{R}$, which is the statement of T15 Remark 4.7 that the Majorana matrix on this route is $v_{R}f$. With $v_{R} \le 4.25 \times 10^{10}$ (the rounding) and $f_{J} \le 4\pi < 4 \cdot 355/113$, $M_{3} < 5.35 \times 10^{11}$ GeV. The conjunction needs $A/(2c) \le T_{\mathrm{desc}}^{4} \le M_{3}^{4}$, so $c \ge c^{*} = A/(2M_{3}^{4})$, and exact evaluation gives $13.49 < c^{*} < 13.50$, more than $5800$ times the per-vector coefficient $2.3 \times 10^{-3}$. Theorem \claimref{synthesis.gauge-algebra-by-tower} gives the odd towers' gauge algebras of dimensions $2, 2, 9, 9$, so their vectors alone give $c = 22 \times 2.3 \times 10^{-3} = 0.0506$; Proposition \claimref{dark-matter.coefficients-floors-they-give-what-gravity}(ii) gives $c \le 0.078$ for a Standard-Model content, and four such towers give $0.312$. At these values $A/(2c) > (4.04\,M_{3})^{4}$ and $> (2.56\,M_{3})^{4}$ respectively, and at $c = 0.312$, $T_{\mathrm{desc}} > 1.36 \times 10^{12}$ GeV. Meeting the signpost on the chain at $c = 0.312$ needs $f_{J} = M_{3}/v_{R} \ge T_{\mathrm{desc}}/v_{R} > 32$, and $m_{D}^{2} = m_{3}M_{3} \ge m_{3}T_{\mathrm{desc}} > (8.2\ \mathrm{GeV})^{2}$ with $m_{3} \ge 49.4$ meV, while $f_{J} \le 4\pi$ gives $m_{D}^{2} = m_{3}f_{J}v_{R} < (5.2\ \mathrm{GeV})^{2}$. Since the comparison is linear in $v_{R}^{4}$, the conclusion at $c = 0.312$ survives any $v_{R}$ below $2.56$ times the stored value; the same chain at one loop has $M_{I} = 3.2 \times 10^{11}$ GeV (Proposition \claimref{standard-model.chain}(2)), about seventeen times the two-loop value, so (i) is a two-loop record-level statement.

\emph{(ii)} Proposition \claimref{dark-matter.ceiling-valued-inherited-success} at its bottom end gives $M_{3} = m_{b}^{2}/m_{3}$; with $m_{b} \le 4.23$ GeV and $m_{3} \ge 49.4$ meV, $M_{3} < 3.63 \times 10^{11}$ GeV and $c^{*} = A/(2M_{3}^{4}) \in (63.7, 63.9)$.

\emph{(iii)} At the top end, $m_{t} \ge 172.28$ GeV and $m_{3} \le 50.5$ meV give $M_{3} > 5.8 \times 10^{14}$ GeV and $(M_{3}/66)^{4} > A/c$ already at the vectors' $c = 0.0506$ without the factor $2$, so the signpost is met; on the chain, $f_{J} = M_{3}/v_{R} > 4000\pi$. Read at the ceiling, with both ends of (U) lowered by a factor $10$, $M_{3}/10 > 5.8 \times 10^{13}$ GeV still satisfies $(M_{3}/10)^{4} > A/c$ at $c = 0.0506$, but $(M_{3}/660)^{4} < A/(2c)$ there (using $M_{3} < 6.1 \times 10^{14}$ GeV from $m_{t} \le 172.86$ GeV and $m_{3} \ge 49.4$ meV), so the clause with the factor $66$ holds only at the common low scale.

The check evaluates every inequality above in exact rational arithmetic, comparing fourth powers so that no root is taken, and recovers $\kappa_{s}$ from its formula and the stored $1.44$ GeV of the chain's Dirac mass at $f_{J} = 1$.
